\label{anlage:skalierung}

Diese Anlage fasst die Rechercheergebnisse zur Skalierung der Wiederverwendung in einer Systemkarte zusammen. Zwei ergänzende Prüftabellen stehen als eigene Anhänge: die dokumentierte Funktionsabdeckung von elf Plattformen (\appendixref{anlage:nachweismatrix}) und deren Integrationsreife als anbindbare Datenquelle (\appendixref{anlage:integrationsreife}).

Die Angaben der Karte wurden aus öffentlich zugänglichen Quellen und den im Bericht genannten Fällen zusammengeführt (Stand August~2026). Sie liegen damit teils nach dem Berichtszeitraum, ergänzen den Bericht und sind nicht Teil der im Zeitraum erbrachten Leistung.

\appendixsection[SM]{Systemkarte}

\begin{landscape}

\begin{Figure}[title=Systemkarte Skalierung der Wiederverwendung, break=false]
\begin{tikzpicture}[semio, x=1mm, y=1mm,
  box/.style={semio cell, inner xsep=0.3em, inner ysep=0.15em},
  band/.style={semio band},
  proc/.style={box, fill=semio-chrome-window, text=semio-chrome-foreground,
               anchor=north west, align=left,
               font=\SemioSans\fontsize{7.2pt}{8.2pt}\selectfont},
  span/.style={draw=semio-chrome-border-normal, fill=semio-chrome-panel, line width=0.75pt},
  guide/.style={draw=semio-chrome-border-normal, line width=0.75pt, opacity=0.28},
  L/.style={font=\SemioSans\fontsize{7.2pt}{8.4pt}\selectfont,
            text=semio-chrome-foreground, anchor=west, align=left},
  Lm/.style={font=\SemioSans\fontsize{7.2pt}{8.4pt}\selectfont,
             text=semio-chrome-text-normal, anchor=west, align=left},
  T/.style={font=\SemioMono\fontsize{7.2pt}{8.4pt}\selectfont,
            text=semio-chrome-foreground, anchor=north west},
  hchip/.style={semio chip, anchor=center,
                font=\SemioMono\fontsize{7.2pt}{8.4pt}\selectfont,
                inner xsep=0.25em, inner ysep=0.1em},
  val/.style={font=\SemioMono\fontsize{7.2pt}{8.4pt}\selectfont,
              text=semio-chrome-foreground, anchor=east},
]
\def\PX{28}
\def\PW{16.6}

%% ============ SYNTHESE-BANNER ==========================================
\node[draw=semio-chrome-border-emphasized, line width=1.5pt, fill=none,
      anchor=north west, text width=173mm, align=left, inner xsep=1.6mm, inner ysep=0.9mm,
      font=\SemioSans\fontsize{7.2pt}{8.5pt}\selectfont, text=semio-chrome-foreground] at (0,131) {%
  {\SemioMono SYNTHESE (E2\SemioSlash{}E3):}\ \ Skalierung entsteht nicht durch eine digitale Plattform allein, sondern aus dem Zusammenspiel von Wiederverwendungsprozessen~(B), digitalen Plattformfunktionen~(C), operativen Diensten~(D) und institutionellen Rahmenbedingungen~(E). Hemmnisse~(F) wirken an den Schnittstellen dieser Ebenen und bestimmen, ob Materialflüsse, Akteursnetzwerke und Prozesse dauerhaft skalieren.};

%% ============ A  SKALIERUNG =============================================
\node[T] at (0,121) {A \ SKALIERUNG \ \textperiodcentered\ \ sechs Dimensionen (E3)};
\foreach \i/\n in {0/Geografische Reichweite, 1/Organisatorische Reichweite,
                   2/Operative Aktivität, 3/Institutionelle Reife,
                   4/Kontinuität, 5/Funktionale Integration}{
  \node[proc, text width=25mm, minimum height=5.5mm] at (\i*29.4,115) {\n};
}

%% ============ G  EVIDENZPFADE ===========================================
\node[T] at (0,109) {G \ EVIDENZPFADE \ \textperiodcentered\ \ durchgezogen quellenbelegt, gestrichelt analytische Zuordnung (E3)};
\foreach \r/\src/\a/\b/\c/\d in {%
  0/EU26/Öffentliche Beschaffung/planbare Nachfrage/Marktbildung/Operative Aktivität,
  1/CCRI25/Regionales Ressourcenkataster/interkommunaler Austausch/größerer Sekundärmarkt/Geografische Reichweite,
  2/UPP25/Förderung und Lagerung/stetigere Materialflüsse/neue Akteure/Institutionelle Reife,
  3/CCRI25/Risikoteilung/Vertrauen/regelmäßige Rückgewinnung/Operative Aktivität}{
  \pgfmathsetmacro\yy{102.5-\r*4.3}
  \node[hchip] at (6,\yy) {\src};
  \node[proc, text width=31mm, minimum height=3.9mm, anchor=west] at (13,\yy) {\a};
  \node[proc, text width=31mm, minimum height=3.9mm, anchor=west] at (49,\yy) {\b};
  \node[proc, text width=37mm, minimum height=3.9mm, anchor=west] at (85,\yy) {\c};
  \node[box, fill=semio-chrome-window, text width=33mm, minimum height=3.9mm, anchor=west,
        font=\SemioSans\fontsize{7.2pt}{8.2pt}\selectfont, text=semio-chrome-foreground] at (127,\yy) {\d};
  \draw[semio edge] (44.2,\yy) -- (48.2,\yy);
  \draw[semio edge] (80.2,\yy) -- (84.2,\yy);
  \draw[semio edge synth] (122.2,\yy) -- (126.2,\yy);
}
\node[Lm] at (163,102.5) {$\uparrow$ auf A};

%% ============ B  D2R PROZESSE ===========================================
\node[T] at (0,85.5) {B \ D2R-PROZESSE \ \textperiodcentered\ \ buildingSMART, neun Prozesse, nicht-lineare Menge};
\node[Lm, text width=26mm] at (0,77.5) {Prozess-Rückgrat};
\foreach \i/\n in {0/Pre-decon\-struction, 1/Decon\-struction, 2/Material Bank Assessment,
                   3/Reuse, 4/Remanu\-facturing, 5/Inventory Manage\-ment,
                   6/Pre-con\-struction Design for Reuse,
                   7/Pre-con\-struction Design with Reusables, 8/Building Design}{
  \pgfmathsetmacro\xx{\PX+\i*\PW}
  \node[proc, text width=14.2mm, minimum height=11mm] at (\xx,82) {\n};
}

%% ---- column guides (span the C and D span zones, behind the spans) -----
\foreach \i in {0,1,2,3,4,5,6,7,8,9}{
  \pgfmathsetmacro\gx{\PX+\i*\PW}
  \draw[guide] (\gx,64.5) -- (\gx,38);
  \draw[guide] (\gx,30.5) -- (\gx,8);
}
\foreach \x/\h in {2/H-05, 4/H-09, 7/H-10}{
  \pgfmathsetmacro\xx{\PX+\x*\PW+\PW/2}
  \node[hchip] at (\xx,69) {\h};
}

%% ============ C  DIGITALE PLATTFORMFUNKTIONEN ===========================
\node[T] at (0,67) {C \ DIGITALE PLATTFORMFUNKTIONEN \ \textperiodcentered\ \ neun Gruppen als Spannen (E3-Zuordnung)};
\foreach \r/\a/\b/\n in {%
  0/0/8/Interoperabilität, 1/0/8/Lebenszyklus, 2/0/2/Bestandserfassung,
  3/2/7/Suche und Matching, 4/6/8/Bedarfserfassung, 5/3/5/Austauschwege,
  6/3/5/Reservierung, 7/2/5/Portfolio, 8/2/8/Auswertung}{
  \pgfmathsetmacro\yy{62-\r*2.85}
  \pgfmathsetmacro\xa{\PX+\a*\PW}
  \pgfmathsetmacro\xb{\PX+\b*\PW+\PW-1.2}
  \node[Lm] at (0,\yy) {\n};
  \draw[span] (\xa,\yy-1.0) rectangle (\xb,\yy+1.0);
}
\foreach \x/\h in {1/H-11, 3/H-07, 5/H-04}{
  \pgfmathsetmacro\xx{\PX+\x*\PW+\PW/2}
  \node[hchip] at (\xx,36.5) {\h};
}

%% ============ D  OPERATIVE DIENSTE ======================================
\node[T] at (0,34.5) {D \ OPERATIVE DIENSTE \ \textperiodcentered\ \ physische Ausführung; ein digitales Match ist noch keine Wiederverwendung};
\foreach \r/\a/\b/\n in {%
  0/0/8/Dienstleister, 1/0/1/Rückbau, 2/1/2/Prüfung, 3/2/3/Zertifizierung,
  4/2/4/Aufbereitung, 5/3/4/Reparatur, 6/2/5/Lagerung, 7/1/5/Logistik}{
  \pgfmathsetmacro\yy{29.5-\r*2.85}
  \pgfmathsetmacro\xa{\PX+\a*\PW}
  \pgfmathsetmacro\xb{\PX+\b*\PW+\PW-1.2}
  \node[Lm] at (0,\yy) {\n};
  \draw[span] (\xa,\yy-1.0) rectangle (\xb,\yy+1.0);
}
\foreach \x/\h in {0/H-01, 3/H-13, 6/H-03, 8/H-12}{
  \pgfmathsetmacro\xx{\PX+\x*\PW+\PW/2}
  \node[hchip] at (\xx,6.5) {\h};
}

%% ============ E  RAHMENBEDINGUNGEN ======================================
\node[T] at (0,4.5) {E \ ÖKOSYSTEMISCHE RAHMENBEDINGUNGEN \ \textperiodcentered\ \ Umfeld des Systems, keine Plattformfunktionen};
\foreach \i/\n in {0/Öffentliche Beschaffung, 1/Normen, 2/Finan\-zierung, 3/Geneh\-migung,
                   4/Fläche und Raum, 5/Risiko und Haftung, 6/Kompe\-tenzen,
                   7/Governance, 8/Ressourcen\-kataster, 9/Wertmodell}{
  \pgfmathsetmacro\xx{\i*17.7}
  \node[proc, text width=15.2mm, minimum height=7.5mm] at (\xx,0.5) {\n};
}

%% ============ H  INSETS (right column) ==================================
\draw[band] (180,-8) rectangle (246,130.5);

%% -- H3  Marktsnapshot 2022 (neben C)
\node[T] at (181.5,130) {H3 \ MARKTSNAPSHOT 2022};
\node[Lm, text width=63mm, anchor=north west] at (181.5,125) {n\,=\,746 Digital Reuse Market Solutions (Websites), Europa, Mai 2022};
\foreach \r/\v/\n in {0/46.4/RBC-Händler, 1/17.2/BC-Händler, 2/15.6/Abbruchunternehmen,
                      3/5.9/Marktplatz, 4/5.6/Lokales Handwerk, 5/4.6/Netzwerkplattform,
                      6/1.9/RBC-Dienstleistung, 7/1.6/RBC-Wissen, 8/1.3/Datenplattform}{
  \pgfmathsetmacro\yy{113-\r*3.15}
  \node[Lm] at (181.5,\yy) {\n};
  \draw[span] (214,\yy-1.05) rectangle (214+\v*0.40,\yy+1.05);
  \node[val] at (243,\yy) {\v};
}
\node[Lm, text width=63mm, anchor=north west] at (181.5,83) {Online-Zahlung: 13,4\,\% aller erfassten Lösungen};

%% -- H2  Hemmnis-Kopplung (neben Reibungszone)
\node[T] at (181.5,77) {H2 \ HEMMNIS-KOPPLUNG};
\node[Lm, text width=63mm, anchor=north west] at (181.5,72) {C-Index (Jaccard-Kookkurrenz), signifikante Paare laut Rakhshan et al.\ 2020, n\,=\,51};
\foreach \c/\n in {0/Kost, 1/Mrkt, 2/Regl, 3/Wahr}{
  \node[val, anchor=south] at (207+\c*9.5,60) {\n};
}
\foreach \r/\n in {0/Markt, 1/Regulierung, 2/Wahrnehmung, 3/Risiko, 4/Gesundheit}{
  \pgfmathsetmacro\yy{57-\r*3.9}
  \node[Lm] at (181.5,\yy) {\n};
}
\foreach \r/\c/\v in {0/0/.44,
                      1/0/.49, 1/1/.45,
                      2/0/.36, 2/1/.35, 2/2/.40,
                      3/0/.34, 3/1/.38, 3/2/.38, 3/3/.63,
                      4/0/.35}{
  \pgfmathsetmacro\yy{57-\r*3.9}
  \pgfmathsetmacro\xx{207+\c*9.5}
  \draw[span] (\xx-4.2,\yy-1.55) rectangle (\xx+4.2,\yy+1.55);
  \node[val, anchor=center] at (\xx,\yy) {\v};
}
\node[Lm, text width=63mm, anchor=north west] at (181.5,38) {Zusammenhang, keine Kausalität (eigene Einordnung, nicht des Papers). Grad eigene Ableitung: Kosten~5, Markt\SemioSlash{}Regulierung\SemioSlash{}Wahrnehmung\SemioSlash{}Risiko je~4, Gesundheit~1.};

%% -- H1  Unterstützung (neben E)
\node[T] at (181.5,25) {H1 \ UNTERSTÜTZUNG FÜR SKALIERUNG};
\node[Lm, text width=63mm, anchor=north west] at (181.5,20) {EU-Umfrage, 19 Antworten aus 10 Ländern; Nenner n\,=\,18 abgeleitet};
\foreach \r/\v/\n in {0/83/Öffentliche Beschaffung, 1/72/Finanzielle Anreize, 2/56/Technische Normen,
                      3/44/Vereinfachte Genehmigung, 4/44/EU-Finanzierung, 5/28/Qualifizierung,
                      6/11/Matching-Plattformen}{
  \pgfmathsetmacro\yy{12-\r*3.0}
  \node[Lm] at (181.5,\yy) {\n};
  \draw[span] (217,\yy-1.05) rectangle (217+\v*0.20,\yy+1.05);
  \node[val] at (243,\yy) {\v\,\%};
}

\end{tikzpicture}
\end{Figure}

\end{landscape}

\clearpage

\appendixsection[LS]{Leseschlüssel}

\SemioTableLong{Ebenen der Systemkarte}{0.16,0.84}{Ebene & Bedeutung}{%
  \SemioTableRow{A & Skalierung: sechs analytische Dimensionen, auf die das System wirkt. Kein Einzelwert; die Dimensionen werden getrennt erhoben.}
  \SemioTableRow{G & Evidenzpfade: quellenbelegte Wirkungsketten mit Quellenchip; der letzte Schritt ist die analytische Zuordnung zu einer Dimension aus A.}
  \SemioTableRow{B & D2R-Prozesse (buildingSMART): neun Prozesse als gemeinsames Prozess-Rückgrat. Nicht-lineare Menge, keine Zeitachse.}
  \SemioTableRow{C & Digitale Plattformfunktionen: neun Gruppen, dargestellt als Spannen über die Prozesse, die sie unterstützen.}
  \SemioTableRow{D & Operative Dienste: physische und professionelle Leistungen, ohne die ein digitales Match nicht zur Wiederverwendung führt.}
  \SemioTableRow{E & Ökosystemische Rahmenbedingungen: institutionelles und marktliches Umfeld, ausdrücklich keine Plattformfunktionen.}
  \SemioTableRow{F & Reibungspunkte: H-Kennungen an den Schnittstellen zwischen den Ebenen.}
  \SemioTableRow{H1--H3 & Empirische Insets neben der Ebene, die sie stützen.}
}

\SemioTableLong{Zeichen und Evidenzgrade}{0.16,0.84}{Zeichen & Bedeutung}{%
  \SemioTableRow{durchgezogener Pfeil & Quellenbelegte Beziehung (E1 direkt berichtet, E2 quellenübergreifend synthetisiert).}
  \SemioTableRow{gestrichelter Pfeil & Analytische Zuordnung dieses Projekts (E3) zu einer Skalierungsdimension; keine Quellenaussage.}
  \SemioTableRow{Balken (Spanne) & Analytische Zuordnung (E3) einer Funktion oder Leistung zu den Prozessen, die sie unterstützt.}
  \SemioTableRow{H-xx & Reibungspunkt an einer Schnittstelle. Die Platzierung ist eine analytische Zuordnung (E3); das Register führt nur die Bezeichnungen, keine Gewichtung.}
  \SemioTableRow{EU26 & Big Buyers Working Together (2026): Public Procurement of Circular Buildings.}
  \SemioTableRow{CCRI25 & CCRI TWG C-BUILD (2025): Final Report on Construction and Circular Buildings.}
  \SemioTableRow{UPP25 & CCRI (2025): Fallbericht Bauteilmarktplatz Uppsala.}
}

\SemioTableLong{Hemmnis-Register}{0.16,0.84}{Kennung & Hemmnis}{%
  \SemioTableRow{H-01 & Kosten}
  \SemioTableRow{H-02 & Akzeptanz, Motivation und Interesse}
  \SemioTableRow{H-03 & Ablauf und Koordination von Rück- und Wiedereinbau}
  \SemioTableRow{H-04 & Logistik}
  \SemioTableRow{H-05 & Informationen über Quellobjekte}
  \SemioTableRow{H-06 & Kompetenzen}
  \SemioTableRow{H-07 & Haftung, Garantie und Versicherung}
  \SemioTableRow{H-08 & Branchenentwicklung}
  \SemioTableRow{H-09 & Technische Qualität der Quellobjekte}
  \SemioTableRow{H-10 & Gesetze, Normen und Vorschriften}
  \SemioTableRow{H-11 & Technische Verfahren für Rück- und Wiedereinbau}
  \SemioTableRow{H-12 & Zusammenarbeit von Akteur:innen}
  \SemioTableRow{H-13 & Beschaffung}
}

\SemioTableLong{Absorbierte Figuren und verbleibende Tabellen}{0.42,0.58}{Absorbiert in die Systemkarte & Verbleibt als Prüftabelle}{%
  \SemioTableRow{D1 Skalierungsdimensionen & D7 Nachweismatrix (18 Funktionen \texttimes\ 11 Plattformen)}
  \SemioTableRow{D2 Drei-Schichten-Architektur & D9 Integrationsreife (Plattform-Bewertung)}
  \SemioTableRow{D3 Plattformfunktionen & Hemmnis-Dossiers und Zuordnungstabelle}
  \SemioTableRow{D8 Mechanismen (Panel B) & Quellen- und Claim-Ledger}
  \SemioTableRow{D8 Umfrage wird zu Inset H1 & C-Index-Volltabelle mit p-Werten}
}

\noindent{\SemioSans\fontsize{7.2pt}{9pt}\selectfont\color{semio-chrome-text-normal}%
Das Register folgt \appendixref{anlage:huerden} und führt die Hemmnisse ausdrücklich ohne Gewichtung. In der Systemkarte erscheinen nur die Kennungen; die Zuordnung einer Kennung zu einer Schnittstelle ist eine analytische Festlegung dieses Projekts (E3) und nicht aus den Quellen abgeleitet. H-02, H-06 und H-08 wirken ebenenübergreifend und sind daher keiner einzelnen Schnittstelle zugeordnet.\par
\vspace{1mm}%
Quellen der Insets: H1 EU-Umfrage (EU26). H2 Rakhshan et al.\ 2020~\cite{rakhshan-components-reuse-review-2020}, C-Index als Jaccard-Kookkurrenz, signifikante Paare laut Tabelle~4, n\,=\,51 ausgewertete Artikel; die Tabelle listet 16 von 190 möglichen Paaren, also nicht alle bestehenden Zusammenhänge. H3 Sivers, Fröhlich und Fivet 2022~\cite{sivers-froehlich-fivet-digital-market-solutions-2022}; Momentaufnahme Mai 2022, kein aktueller Marktanteil.\par}
